\clearpage

% --------------------------------------------------------------------------------------------

\section {Summary}
\label{sec:summary}

In this technical note,
we have presented the results of the 3-flavour $\bar{\nu}/\nu$ joint analysis
performed by the {\em VALOR} group,
on the combined Run 1-6 data set,
for an exposure of \TotalPotNu{} POT in neutrino mode
and \TotalPotNuBar{} POT in antineutrino mode. 
The goal of the analysis is to search for CP violation in the leptonic
and put limits on $\delta_{CP}$ as well as 
measure all the other oscillation parameters to which T2K is sensitive:
$\sin^2 \theta_{13}$, $\sin^2 \theta_{23}$, $\Delta m^2_{32}$ ($\Delta m^2_{13}$)
and partially mass hierarchy, to which T2K has little sensitivity.

This analysis predicts 
\NskOscExpectationNumu{} $\pm$ \NskOscExpectationNumuError{} (syst) for neutrino $\mu$-like,
\NskOscExpectationNumuBar{} $\pm$ \NskOscExpectationNumuBarError{} (syst) for antineutrino $\mu$-like,
\NskOscExpectationNue{} $\pm$ \NskOscExpectationNueError{} (syst) for neutrino $e$-like,
and
\NskOscExpectationNueBar{} $\pm$ \NskOscExpectationNueBarError{} (syst) for antineutrino $e$-like
single ring events in Super K
and observes respectively 
\NskObservedNumu{} for neutrino $\mu$-like,
\NskObservedNumuBar{}  for antineutrino $\mu$-like,
\NskObservedNue{}  for neutrino $e$-like,
\NskObservedNueBar{}  for antineutrino $e$-like.

The $\bar{\nu}/\nu$ analysis was performed
in a framework of 3-flavour oscillations including matter effects in constant-density matter.
The observed reconstructed energy spectrum of single $\mu$-like ring events
and reconstructed energy versus lepton angle of single $e$-like ring events
was used to compute the confidence intervals.
All \NSyst{} systematic parameters as well as the oscillation parameters
which were not parameters of interest considered in this analysis were marginalised.

From the fit of the Run 1-6 data we can see that the favored value of $\delta_{CP}$ is -1.854, close to $-\pi/2$.
When the constant $\Delta \chi^2$ method is used $\delta_{CP}=0$ is excluded at 90\% CL
and a large portion between $0$ and $+\pi$ is excluded at $2\sigma$ if normal hierarchy is considered.
In case of inverted hierarchy a not negligible range of $\delta_{CP}$ around $+\pi/2$ is excluded with 
a significance of $3\sigma$.

From the contour of $\sin^2 \theta_{13}$ Vs $\delta_{CP}$ without the reactor constraint 
it is clear that the T2K data are in agreement with the reactor measurements within $1\sigma$ 
if $\delta_{CP} \sim -\pi/2$.

It has also been shown that the result is dominated by the excess of event in the neutrino mode
$e$-like sample, though the antineutrino mode $e$-like also prefers values of $\delta_{CP}$ 
close to $-\pi/2$, though the significance is not high.


%\vspace{3mm}
%{\huge{
%TODO: ADD COMMENTS ABOUT BEST-FIT OSCILLATION PARAMETERS
%AND 1SIGMA RANGES
%}}
%\vspace{3mm}

%\vspace{3mm}
%{\huge{
%TODO: ADD COMMENTS ABOUT GOODNESS OF FIT
%}}
%\vspace{3mm}

%\vspace{3mm}
%{\huge{
%TODO: ADD COMMENTS ABOUT 1D CONTOURS
%}}
%\vspace{3mm}
%
%\vspace{3mm}
%{\huge{
%TODO: ADD COMMENTS ABOUT 2D CONTOURS
%}}
%\vspace{3mm}





\clearpage
